% ── sections/03_theory.tex ───────────────────────────────────────────────────

\section{Theory and Model Architecture}
\label{sec:theory}

\subsection{Per-Network Graph Neural Network Encoder}

Let $\mathcal{G}^{(k)} = (\mathcal{V}^{(k)}, \mathcal{E}^{(k)})$ denote the $k$-th PPI
network with node feature matrix $\mathbf{X}^{(k)} \in \mathbb{R}^{N_k \times D}$, where
$N_k$ is the number of genes in network $k$ and $D = 64$ is the multi-omics feature
dimension. A shared $L$-layer GNN encoder produces node embeddings:
\begin{equation}
  \mathbf{H}^{(k,l)} = \sigma\!\left(
    \text{GNN}_l\!\left(\mathbf{H}^{(k,l-1)},\, \mathbf{A}^{(k)}\right)
  \right), \quad l = 1, \ldots, L,
  \label{eq:gnn_layer}
\end{equation}
where $\mathbf{H}^{(k,0)} = \mathbf{X}^{(k)}$, $\mathbf{A}^{(k)}$ is the symmetrically
normalised adjacency matrix with added self-loops, and $\sigma$ is LeakyReLU. We support
four GNN variants:

\textbf{GCN}~\cite{gcn2017}: $\text{GNN}(\mathbf{H},\mathbf{A}) = \hat{\mathbf{A}}\mathbf{H}\mathbf{W}$,
where $\hat{\mathbf{A}} = \tilde{\mathbf{D}}^{-1/2}\tilde{\mathbf{A}}\tilde{\mathbf{D}}^{-1/2}$
is the symmetric normalised adjacency.

\textbf{GAT}~\cite{gat2018}: Attention coefficients $\alpha_{ij}$ are learned via a shared
attentional mechanism over node pairs, allowing the model to weight neighbour contributions
adaptively.

\textbf{GIN}~\cite{gin2019}: $\text{GNN}(\mathbf{H},\mathbf{A}) =
\text{MLP}\!\left((1+\epsilon)\mathbf{H} + \mathbf{A}\mathbf{H}\right)$, maximally
expressive within the Weisfeiler--Lehman hierarchy.

\textbf{GraphSAGE}~\cite{sage2017}: $\text{GNN}(\mathbf{H},\mathbf{A}) =
\mathbf{W}\,[\mathbf{h}_u \,\|\, \text{MEAN}_{v \in \mathcal{N}(u)}\mathbf{h}_v]$,
concatenating a node's own embedding with its mean neighbourhood embedding.

When residual connections are enabled, layers $l \ge 2$ add a skip connection:
\begin{equation}
  \mathbf{H}^{(k,l)} = \sigma\!\left(\text{GNN}_l\!\left(\mathbf{H}^{(k,l-1)},
  \mathbf{A}^{(k)}\right)\right) + \mathbf{H}^{(k,l-1)}.
  \label{eq:residual}
\end{equation}

\subsection{Meta-Graph Construction}

After per-network encoding, a meta-graph $\mathcal{G}^{\text{meta}}$ is constructed.
Each unique gene across all $K$ networks receives one meta-node; directed edges connect
each per-network copy of gene $i$ (in network $k$) to its meta-node. Meta-node features
$\mathbf{X}^{\text{meta}} \in \mathbb{R}^{M \times D}$ are initialised by assigning each
meta-node the multi-omics features of the corresponding gene (from whichever network first
provides them), where $M$ is the number of unique genes. This design aggregates
cross-network evidence without redundancy and scales linearly with the number of unique genes.

\subsection{Learnable Per-Network Importance Weights}

In the standard EMGNN, all network encoders contribute equally to the meta-graph.
We introduce a trainable weight vector $\boldsymbol{\theta} \in \mathbb{R}^{K}$, producing
normalised network importances:
\begin{equation}
  \mathbf{w} = \text{softmax}(\boldsymbol{\theta}), \quad
  w_k = \frac{e^{\theta_k}}{\sum_{j=1}^{K} e^{\theta_j}}.
  \label{eq:network_weights}
\end{equation}
The embedding of each input node in network $k$ is scaled by $w_k$ before being
passed to the meta-graph GNN, allowing high-quality networks to be up-weighted
end-to-end during training.

\subsection{Meta-Graph GNN and Classifier}

A second $L$-layer GNN (same backbone) operates on $\mathcal{G}^{\text{meta}}$,
aggregating cross-network signals over meta-node edges:
\begin{equation}
  \mathbf{Z} = \text{GNN}^{\text{meta}}\!\left(\mathbf{X}^{\text{meta}},\,
  \mathbf{A}^{\text{meta}}\right) \in \mathbb{R}^{M \times H}.
  \label{eq:meta_gnn}
\end{equation}
Each meta-node embedding $\mathbf{z}_i$ is passed through a two-layer MLP with softmax
output to produce cancer-class probabilities $\hat{p}_i$. The model is trained with
cross-entropy loss, optionally with label smoothing as in Eq.~\ref{eq:labelsmoothing}.

\subsection{Normalisation in Full-Batch Graph Learning}

Standard Batch Normalisation~\cite{batchnorm2015} normalises over the batch dimension:
\begin{equation}
  \hat{h}_i = \frac{h_i - \mu_\mathcal{B}}{\sqrt{\sigma^2_\mathcal{B} + \epsilon}},
  \quad \mu_\mathcal{B} = \frac{1}{N}\sum_{i=1}^N h_i,
  \label{eq:batchnorm}
\end{equation}
where $\mathcal{B}$ is the set of all $N$ nodes. In the full-batch graph setting there is
no variation across batches, so the running statistics computed during training equal those
at inference; the normalisation is effectively a fixed affine transform and provides no
regularisation benefit.

GraphNorm~\cite{graphnorm2021} instead computes per-graph statistics, subtracting a
learnable mean shift $\alpha \cdot \mu_{\mathcal{G}}$ before normalisation:
\begin{equation}
  \hat{h}_i = \frac{h_i - \alpha\,\mu_{\mathcal{G}}}
              {\sqrt{\sigma^2_{\mathcal{G}} + \epsilon}},
  \quad \mu_{\mathcal{G}} = \frac{1}{|\mathcal{V}|}\sum_{i \in \mathcal{V}} h_i,
  \label{eq:graphnorm}
\end{equation}
where $\alpha \in [0,1]$ is a trainable scalar that controls the degree of mean removal.
This preserves meaningful within-graph distributional information that \texttt{BatchNorm1d}
discards in the full-batch regime.

Our improved model supports \texttt{norm\_type} $\in$ \{\texttt{batch}, \texttt{graph},
\texttt{layer}, \texttt{none}\}, implemented via \texttt{BatchNorm1d}, \texttt{GraphNorm}
(from PyTorch Geometric), \texttt{LayerNorm}, and identity respectively. Ablation results
in Section~\ref{sec:ablation} demonstrate the cost of \texttt{batch} normalisation; direct
comparison of the remaining options is future work.

\subsection{Integrated Gradients for Interpretability}

Integrated Gradients~\cite{ig2017} satisfies the \emph{completeness} axiom: the sum of
all feature attributions equals the difference between the model output at the input and
at the baseline. For a gene with feature vector $\mathbf{x}$ and zero baseline
$\mathbf{x}' = \mathbf{0}$:
\begin{equation}
  \text{IG}_d(\mathbf{x}) =
    x_d \int_0^1
    \frac{\partial F(\alpha\mathbf{x})}{\partial x_d} \, d\alpha,
  \label{eq:ig}
\end{equation}
where $F$ is the cancer-class logit and the integral is approximated with 50 Riemann
steps. This property ensures that attributions account for the full prediction score, making
Integrated Gradients preferable to vanilla gradient methods for model-level interpretability.

\textbf{Baseline choice.} We use the zero vector $\mathbf{x}' = \mathbf{0}$ as the
reference point, which is the most common choice in the IG literature. This corresponds to
asking: ``how much does each feature contribute relative to a gene with all omics
measurements set to zero?'' The zero baseline is interpretable when features are
pre-normalised (as in the EMOGI/EMGNN pipeline), since zero approximately represents the
absence of signal or a population-average reference. A potential limitation is that the
zero vector may not lie in the natural data manifold; alternative baselines such as a
Gaussian blur or the mean across non-cancer genes could be explored in future work to
assess baseline sensitivity. Attribution magnitudes reported in this work should therefore
be interpreted as importance \emph{relative to the zero baseline}, not as absolute measures
of biological relevance.

